\documentclass[11pt]{article}

% ── Packages ─────────────────────────────────────────────────────────────────
\usepackage[utf8]{inputenc}
\usepackage[T1]{fontenc}
\usepackage[margin=1in]{geometry}
\usepackage{hyperref}
\usepackage{url}
\usepackage{booktabs}
\usepackage{xcolor}
\usepackage[protrusion=true,expansion=false]{microtype}
\usepackage{fancyhdr}
\usepackage{titlesec}
\usepackage{listings}
\usepackage{tcolorbox}
\usepackage{array}
\usepackage{tabularx}
\usepackage{multirow}
\usepackage{natbib}
\usepackage{graphicx}
\usepackage{amsmath}
\usepackage{amssymb}
\usepackage{setspace}
\usepackage{parskip}

% ── Colors ────────────────────────────────────────────────────────────────────
\definecolor{aira-blue}{RGB}{31,78,121}
\definecolor{aira-mid}{RGB}{46,117,182}
\definecolor{aira-light}{RGB}{232,240,248}
\definecolor{aira-warn}{RGB}{255,248,232}
\definecolor{aira-code}{RGB}{244,244,244}
\definecolor{aira-red}{RGB}{192,0,0}

% ── Hyperref setup ────────────────────────────────────────────────────────────
\hypersetup{
  colorlinks=true,
  linkcolor=aira-mid,
  citecolor=aira-mid,
  urlcolor=aira-mid,
  pdftitle={AIRA -- AI-Induced Risk Audit},
  pdfauthor={William M. Parris},
  pdfsubject={Software Engineering, AI Safety},
  pdfkeywords={AI-generated code, fail-soft, RLHF, static analysis, software reliability}
}

% ── Section formatting ────────────────────────────────────────────────────────
\titleformat{\section}{\large\bfseries\color{aira-blue}}{{\thesection}}{1em}{}
\titleformat{\subsection}{\normalsize\bfseries\color{aira-mid}}{{\thesubsection}}{1em}{}
\titleformat{\subsubsection}{\normalsize\bfseries\itshape\color{aira-mid}}{{\thesubsubsection}}{1em}{}

% ── Header / footer ───────────────────────────────────────────────────────────
\pagestyle{fancy}
\fancyhf{}
\fancyhead[L]{\small\textit{AIRA --- AI-Induced Risk Audit}}
\fancyhead[R]{\small\textit{BDB Labs, 2026}}
\fancyfoot[C]{\thepage}
\renewcommand{\headrulewidth}{0.4pt}

% ── Code listing style ────────────────────────────────────────────────────────
\lstset{
  backgroundcolor=\color{aira-code},
  basicstyle=\ttfamily\small,
  breaklines=true,
  frame=single,
  rulecolor=\color{gray!30},
  xleftmargin=1em,
  xrightmargin=1em,
  aboveskip=0.5em,
  belowskip=0.5em
}

% ── tcolorbox styles ─────────────────────────────────────────────────────────
\tcbuselibrary{skins,breakable}

\newenvironment{quotebox}{%
  \begin{tcolorbox}[
    enhanced,
    colback=aira-light,
    colframe=aira-mid,
    leftrule=4pt, rightrule=0pt, toprule=0pt, bottomrule=0pt,
    arc=0pt,
    left=6pt, right=6pt, top=4pt, bottom=4pt
  ]
  \itshape
}{%
  \end{tcolorbox}
}

\newenvironment{definitionbox}{%
  \begin{tcolorbox}[
    enhanced,
    colback=aira-light!60,
    colframe=aira-blue,
    arc=2pt,
    left=8pt, right=8pt, top=6pt, bottom=6pt
  ]
}{%
  \end{tcolorbox}
}

% ── Utility commands ──────────────────────────────────────────────────────────
\newcommand{\checkid}[1]{\texttt{#1}}
\newcommand{\aira}{\textsc{aira}}
\newcommand{\HIGH}{\textcolor{aira-red}{\textbf{HIGH}}}
\newcommand{\MED}{\textcolor{orange}{\textbf{MEDIUM}}}
\newcommand{\LOW}{\textcolor{gray}{\textbf{LOW}}}

% ─────────────────────────────────────────────────────────────────────────────
%  DOCUMENT
% ─────────────────────────────────────────────────────────────────────────────
\begin{document}

\begin{titlepage}
  \vspace*{2cm}
  {\Huge\bfseries\color{aira-blue} AIRA --- AI-Induced Risk Audit\par}
  \vspace{0.5em}
  {\Large\itshape A Structured Inspection Framework for AI-Generated Code\par}
  \vspace{2em}
  {\large \textbf{William M.~Parris}\par}
  \vspace{0.4em}
  {\normalsize BDB Labs $\mid$ Jurisprudential AI Governance Initiative\par}
  {\normalsize \href{mailto:bill@bageltech.net}{bill@bageltech.net} $\quad$ \href{https://aira.bageltech.net}{aira.bageltech.net}\par}
  \vspace{0.4em}
  {\normalsize 2026\par}
  \vspace{2em}
  \noindent\rule{\textwidth}{0.4pt}
\end{titlepage}

% ── Abstract ──────────────────────────────────────────────────────────────────
\begin{abstract}
Engineers who work extensively with AI coding tools observe a directional pattern:
AI-generated code often fails quietly, preserving the surface appearance of function
while degrading or concealing guarantees. This paper introduces the
\emph{Reward-Shaped Failure Hypothesis} --- the proposal that this pattern may be a structural
artifact of how AI systems are optimized through human feedback, rather than merely a random
distribution of bugs. Systems trained to maximize positive evaluation signals may learn to
suppress visible failure, because a crashing program is penalized more heavily than a
program that returns a wrong or degraded result. We define \emph{failure truthfulness} as
the property that a system's observable outputs accurately represent its internal success or
failure state. We then present \aira{} (AI-Induced Risk Audit), a 15-check inspection
framework designed to detect the specific class of failure this optimization pressure may
produce. \aira{} is
implemented as a deterministic static analysis engine with optional LLM augmentation,
available as both a CLI tool and a web scanner.

We report empirical results from two studies: a full deterministic audit of a 1,120-file
production Constitutional AI governance system (3,297 findings, 13 of 15 checks failing,
zero passing), and a rebuilt balanced public-corpus pilot of 600 files (300 agent-attributed,
300 matched human controls) showing a 1.32$\times$ overall excess of high-severity findings
in AI-attributed code, with the strongest separation in JavaScript (3.18$\times$) and
the clearest check-level support in exception-handling-related patterns. A secondary
comparison using a cloud LLM evaluator produced findings at a 44:1 ratio below the
deterministic scanner, consistent with the suppression pattern the hypothesis predicts.
\aira{} is designed for governance, compliance, and safety-critical systems where
fail-closed behavior is often a definitional requirement.

\medskip\noindent\textbf{Keywords:} AI-generated code, fail-soft behavior, RLHF,
failure truthfulness, static analysis, software reliability, code audit
\end{abstract}

% ─────────────────────────────────────────────────────────────────────────────
\section{Introduction}
% ─────────────────────────────────────────────────────────────────────────────

The question traditionally asked of software is whether it works. Testing frameworks,
code review practices, and quality metrics are organized around that question. They catch
errors of commission --- code that produces the wrong output, throws an unexpected exception,
or fails to compile. What they are not calibrated to catch is a different class of failure:
code that continues executing, returns a value, and reports success, while silently violating
the guarantees it was supposed to maintain.

This paper argues that AI-assisted development may increase the prevalence and consistency
of exactly this class of failure. The argument is structural rather than anecdotal. AI coding
systems are optimized in part through feedback signals from human evaluators. A program that
crashes produces a strongly negative signal. A program that returns a result --- even a
degraded or incorrect one --- produces a far less penalized signal. Over training, this
asymmetry may shape code generation toward the \emph{appearance} of correctness rather than
correctness under all conditions.

The resulting pattern appears directional rather than random: AI-generated code may
disproportionately fail in ways that preserve surface function while concealing broken
guarantees. Standard review practices, designed for human error patterns, may be
insufficiently calibrated to detect it.

We introduce \aira{} (AI-Induced Risk Audit) as a targeted response to this problem. \aira{}
is not a general-purpose code quality tool. It is a structured inspection framework organized
around a single question:

\begin{quotebox}
Not ``does this code work?'' but ``does this code tell the truth about whether it is working?''
\end{quotebox}

The remainder of this paper is organized as follows. Section~\ref{sec:theory} develops the
Reward-Shaped Failure Hypothesis and formally defines failure truthfulness.
Section~\ref{sec:related} positions \aira{} relative to existing work.
Section~\ref{sec:framework} presents the framework design, architecture, and 15 checks.
Section~\ref{sec:empirical} reports empirical results from two studies.
Section~\ref{sec:output} describes the output specification.
Section~\ref{sec:limitations} addresses limitations.
Section~\ref{sec:conclusion} concludes with future directions.

\subsection{Contributions}

This work makes the following contributions:
\begin{itemize}
  \item Introduces the \emph{Reward-Shaped Failure Hypothesis} --- a structural account of
        why AI-assisted codebases may exhibit systematic fail-soft behavior
  \item Defines \emph{failure truthfulness} as a formal system property distinct from
        correctness
  \item Presents \aira{}, a 15-check inspection framework targeting the specific failure
        class this hypothesis predicts
  \item Reports empirical findings from a full \aira{} audit of a 1,120-file production
        AI governance system
  \item Reports a rebuilt balanced 600-file corpus pilot (300 AI-attributed vs.\ 300
        human-authored) showing a 1.32$\times$ overall and 3.18$\times$ JavaScript
        high-severity differential
  \item Shows that an LLM evaluator applied to the same surfaces exhibits the same
        suppression pattern at a 44:1 ratio below deterministic detection
  \item Releases \aira{} as an open-source tool with CLI, web, and research collection
        components at \url{https://github.com/BDB-Labs/aira-scanner}
\end{itemize}

% ─────────────────────────────────────────────────────────────────────────────
\section{Theoretical Foundation}
\label{sec:theory}
% ─────────────────────────────────────────────────────────────────────────────

\subsection{The Pattern}

Engineers who work extensively with AI coding tools begin to notice a pattern. The code
compiles. The tests pass. The function returns a value. And yet, under failure conditions,
the system does not stop --- it continues, quietly, in a degraded state that nothing in the
output makes visible. Errors are absorbed. Exceptions are swallowed. Guarantees written into
the design simply do not propagate when they are violated.

This is not merely a random distribution of bugs. The observed pattern is directional:
AI-generated code often fails quietly, in ways that preserve the surface appearance of
function. Standard code review practices may not be calibrated to detect this class of
failure, because human developers may not produce it at the same rate or in the same form.

In a coding context, this pattern most commonly manifests as:
\begin{itemize}
  \item \textbf{Silent exception handling} --- because a swallowed error looks like continued
        success
  \item \textbf{Graceful degradation without disclosure} --- because a degraded response feels
        better than a hard stop
  \item \textbf{Optimistic return values} --- because returning \texttt{None} or raising feels
        like giving up
  \item \textbf{Broad try/catch blankets} --- because a crash is the worst possible visible
        outcome
  \item \textbf{Happy-path test coverage} --- because passing tests generate positive feedback
\end{itemize}

\subsection{The Reward-Shaped Failure Hypothesis}

\aira{} offers a hypothesis about the structural cause. AI systems are optimized, in part,
through feedback signals that reward outputs humans evaluate positively. A crashing program
is a strongly negative signal. A program that returns a result --- even a wrong one, even a
degraded one --- is far less penalized. Over training, this pressure shapes how these systems
generate code: not toward correctness under all conditions, but toward the \emph{appearance}
of correctness under the conditions most likely to be evaluated.

\begin{quotebox}
The failure patterns documented in this framework are a predictable consequence of how AI
systems are optimized --- a structural artifact of the feedback environment, not a defect
in any particular implementation.
\end{quotebox}

Because the failure distribution is non-random, it is auditable. The same optimization
pressure that produces silent exception handling in one codebase will produce it in another.
\aira{} is organized around that predictability.

This hypothesis connects two conversations that have largely proceeded in parallel. Software
reliability researchers have documented increasing rates of silent failure in AI-assisted
codebases. AI safety researchers have identified that RLHF inherently trains models toward
fail-open behavior \citep{casper2023open}. \aira{} bridges these observations: the alignment
process itself is the structural cause of the software vulnerability.

\subsection{Failure Truthfulness}
\label{sec:failure-truthfulness}

The innovation in \aira{} is not detection of bugs per se, but detection of a specific
epistemic failure. We define:

\begin{definitionbox}
  \textbf{Failure Truthfulness:} The property that a system's observable outputs accurately
  represent its internal success or failure state, without suppression, ambiguity, or
  degradation masking.
\end{definitionbox}

A system can be functionally correct on the happy path and \emph{failure-untruthful} under
error conditions. Standard correctness testing does not distinguish these. \aira{} is
organized around failure truthfulness as a first-class system property, separate from and
complementary to correctness.

Failure truthfulness fails when a system returns a success signal after a critical operation
has failed (Check C01), when audit evidence can be lost without halting execution (C02), when
exceptions are absorbed without preserving failure semantics (C03), or when outputs are
presented as authoritative without a confidence posture (C13). These are not independent bugs
--- they are different surfaces of the same underlying property violation.

\subsection{Implications for Governance-Critical Systems}

The stakes are highest in governance, compliance, and safety-critical systems --- where
fail-closed behavior is often not a design preference but a definitional requirement. A
system that must refuse to act under uncertainty is difficult to build reliably on code that
was shaped to suppress uncertainty signals.

\aira{} was developed in part through direct observation of this tension during construction
of production Constitutional AI governance frameworks, where AI coding agents consistently
introduced silent exception handlers into audit-critical paths --- not through any individual
mistake, but as a systematic pattern across hundreds of generated functions.

In distributed systems, the failure mode compounds. A swallowed exception in a microservice
does not only affect a local operation --- it can corrupt downstream data pipelines, poison
distributed caches, and propagate a false success signal across system boundaries.

% ─────────────────────────────────────────────────────────────────────────────
\section{Related Work and Positioning}
\label{sec:related}
% ─────────────────────────────────────────────────────────────────────────────

\aira{} occupies a specific position in the landscape of code quality and AI safety tooling.

General-purpose static analysis tools --- SonarQube, Semgrep, CodeClimate, Pylint --- detect
broad classes of defect including security vulnerabilities, style violations, and known bug
patterns. They are not designed to detect the directional failure bias introduced by
optimization pressure. An AI-generated codebase may satisfy standard static analysis while
still exhibiting extensive fail-soft behavior, because the failure mode does not primarily
manifest as syntactic error or known vulnerability pattern --- it manifests as structurally
valid code with the wrong semantics under failure conditions.

Observability and failure transparency research has long recognized the distinction between
fail-open and fail-closed system design \citep{beyer2016sre}. \aira{} formalizes these
concerns as an auditable checklist and connects them to a specific structural cause in
AI-assisted development.

Research on LLM code generation quality has documented that AI-generated code exhibits higher
rates of missing boundary checks, empty exception handlers, and reduced robustness compared
to human-written equivalents. Large-scale structural comparisons of agentic versus
human-authored pull requests using the AIDev dataset have confirmed that AI-generated code is
structurally distinct from human-written code across commit structure and breadth of file
modifications \citep{ogenrwot2026how}. Empirical analysis of test failures in AI-generated
PRs has found that runtime errors dominate at 62.6\% versus 37.4\% compile-time failures,
with assertion failures as the most common issue \citep{msr2026testfailures} --- a
distribution consistent with fail-soft behavior. In AI safety, the fail-open tendency of
RLHF-trained models has been discussed primarily in the context of conversational evasion
\citep{casper2023open}. \aira{} applies this observation to software architecture.

\subsection{Distinction from Prior Work}

While recent studies have documented higher rates of robustness issues and bugs in
LLM-generated code --- such as missing boundary checks, empty exception handlers, and reduced
performance under edge cases --- most focus on measuring bug prevalence or functional
correctness without identifying a unifying structural cause or providing targeted auditing
instrumentation. Benchmarks such as CoderEval \citep{yu2024codereval} have evaluated
pragmatic code generation and shown varying robustness, while systematic surveys catalog
functional, security, and other bug types in AI-generated code \citep{gao2025survey}.

\aira{} advances the literature by (1) naming the Reward-Shaped Failure Hypothesis as the
underlying optimization-driven mechanism, (2) defining failure truthfulness as a distinct,
measurable system property complementary to traditional correctness, and (3) delivering a
deterministic 15-check framework explicitly engineered to detect the predicted failure class
while remaining resistant to the same suppression tendencies. This bridges software
reliability and observability research with AI alignment concerns in a directly actionable
way for governance-critical systems.

% ─────────────────────────────────────────────────────────────────────────────
\section{The AIRA Framework}
\label{sec:framework}
% ─────────────────────────────────────────────────────────────────────────────

\subsection{Design Principles}

\aira{} was designed around four principles that distinguish it from general code quality
tools:

\textbf{Hypothesis-derived checks.} The 15 checks target specific failure modes that the
Reward-Shaped Failure Hypothesis predicts. The framework is organized around a single
question: does this code tell the truth about whether it is working?

\textbf{Deterministic backbone.} The primary scanner is parser-backed deterministic analysis
--- not LLM-assisted classification. This ensures the instrument itself is not susceptible
to the failure mode it is designed to detect. LLM augmentation is available but treated as
optional enrichment, not ground truth.

\textbf{Fail-closed semantics.} \texttt{UNKNOWN} is not a neutral result. Two checks (C07,
C12) require human review. On governance-critical paths, \texttt{UNKNOWN} should be treated
as a conditional \texttt{FAIL} pending manual verification.

\textbf{Research posture.} \aira{} makes no claim that a passing result means a system is
safe. The correct language for \aira{} outputs is \emph{observed}, \emph{measured}, and
\emph{suggests} --- not \emph{proves} or \emph{guarantees}.

\subsection{System Architecture}

\aira{} is implemented as four connected components:

\textbf{Deterministic rule engine.} Parser-backed static analysis for Python (AST-based),
JavaScript, TypeScript, and JSX/TSX (esprima-backed with lexical fallback).

\textbf{CLI tool.} Supports static, LLM, and hybrid scan modes across local repos. Includes
Ollama model discovery, selected-model validation, and aggregate-only research submission
via the research flag (\path{--submit-research-aggregate}). Ollama is treated as a generic
abstraction layer.

\textbf{Web scanner.} Provider-routed scan surface with fallback hierarchy: configured cloud
or Ollama route $\to$ deterministic server-side static scan $\to$ browser heuristics.

\textbf{Research collection pipeline.} Aggregate-only submission to Supabase (hosted) or
JSONL (local/CI). Raw source code, file paths, and snippets are never transmitted.

\subsection{Scan Modes}

\aira{} supports three scan modes. \emph{Static} mode uses deterministic built-in analysis
and is the canonical baseline for research. \emph{LLM} mode uses provider-assisted analysis
and is useful for exploratory comparison. \emph{Hybrid} mode merges static and LLM findings,
allowing provider-assisted signal to supplement the deterministic baseline.

\subsection{The 15 Checks}
\label{sec:checks}

Each check targets a specific failure mode predicted by the Reward-Shaped Failure Hypothesis.
Table~\ref{tab:checks} summarizes the full check set.

\begin{table}[htbp]
\centering
\caption{AIRA Check Summary}
\label{tab:checks}
\small
\begin{tabularx}{\textwidth}{@{}ll>{\raggedright\arraybackslash}p{3.2cm}>{\raggedright\arraybackslash}X@{}}
\toprule
\textbf{ID} & \textbf{Name} & \textbf{Automation} & \textbf{Primary Concern} \\
\midrule
C01 & Success Integrity & Automated & Success returned after critical failure \\
C02 & Audit / Evidence Integrity & Automated & Audit or evidence loss occurs silently \\
C03 & Broad Exception Suppression & Automated & Exceptions swallowed or neutralized \\
C04 & Distributed Fallback & Automated (heuristic) & Fallback scattered, weakens guarantees \\
C05 & Bypass / Override Paths & Automated & Flags or overrides disable safeguards \\
C06 & Ambiguous Return Contracts & Automated & \texttt{None}/\texttt{null} blurs absence and failure \\
C07 & Parallel Logic Drift & Human review only & Divergent paths produce inconsistent semantics \\
C08 & Unsupervised Background Tasks & Automated & Async work lacks supervision \\
C09 & Environment-Dependent Safety & Automated & Safety relaxed in dev/staging paths \\
C10 & Startup Integrity & Automated & Init catches failure and continues \\
C11 & Deterministic Reasoning Drift & Automated & Decision paths use non-deterministic settings \\
C12 & Source-to-Output Lineage & Human review only & Outputs lack traceable origin \\
C13 & Confidence Misrepresentation & Automated & Degraded output presented without confidence posture \\
C14 & Test Coverage Asymmetry & Automated & Failure-path tests lag happy-path tests \\
C15 & Retry / Idempotency Drift & Automated & Retries on writes without idempotency controls \\
\bottomrule
\end{tabularx}
\end{table}

To illustrate the distinction between AI-typical and human-typical failure handling,
consider C01 (Success Integrity):

\begin{lstlisting}[language=Python, caption={AI-typical: failure concealed}]
def process(data):
    try:
        result = persist(data)
        return {"status": "ok"}   # success returned regardless
    except Exception:
        log.warning("persist failed")
        return {"status": "ok"}   # failure concealed
\end{lstlisting}

\begin{lstlisting}[language=Python, caption={Human-typical: failure propagated}]
def process(data):
    result = persist(data)        # raises on failure
    if not result:
        raise PersistenceError("persist returned empty")
    return {"status": "ok"}
\end{lstlisting}

The AI version is not wrong on the happy path. It is \emph{failure-untruthful}: it returns
success after a broken guarantee.

\subsection{Check Co-occurrence Patterns}

The 15 checks are not fully independent. Certain co-occurrence patterns describe
characteristic failure profiles:
\begin{itemize}
  \item \checkid{C03} $+$ \checkid{C01} often indicates explicit failure concealment:
        exceptions are swallowed to preserve a success-signaling return path
  \item \checkid{C04} $+$ \checkid{C09} often indicates environment-shaped degraded
        assurance: fallback behavior is distributed and relaxed in non-production environments
  \item \checkid{C02} $+$ \checkid{C10} often indicates a system that can start and report
        readiness despite losing evidence guarantees at initialization
  \item \checkid{C13} is frequently the epistemic surface that makes the rest look normal:
        confidence misrepresentation makes other failures invisible to downstream consumers
\end{itemize}

Reviewers should treat co-occurring check failures as failure profiles, not independent
findings.

\subsection{Reviewer Instructions}

For each check, reviewers mark:
\begin{itemize}
  \item \texttt{PASS} --- requirement is fully satisfied
  \item \texttt{FAIL} --- requirement is violated (include: file path, line number,
        description)
  \item \texttt{UNKNOWN} --- cannot be determined without runtime or additional context
\end{itemize}

\textbf{On UNKNOWN:} \texttt{UNKNOWN} does not mean safe. It means the scanner cannot
responsibly automate the conclusion. For governance-critical checks, \texttt{UNKNOWN} should
be treated as a conditional \texttt{FAIL} pending manual runtime verification.

% ─────────────────────────────────────────────────────────────────────────────
\section{Empirical Validation}
\label{sec:empirical}
% ─────────────────────────────────────────────────────────────────────────────

This section reports results across two studies. Study~1 is a full deterministic \aira{}
audit of a production Constitutional AI governance system, presented as an existence proof
that the predicted failure class can be real, pervasive, and systematic in practice.
Study~2 is a rebuilt balanced corpus pilot across AI-attributed and human-authored
codebases, designed to test whether the pattern generalizes. Together they move the
Reward-Shaped Failure Hypothesis from structural argument toward a measured, though still
conditional, empirical signal.

\subsection{Study Design}

Both studies use the \aira{} deterministic static engine exclusively. LLM-assisted scanning
is reported separately in Section~\ref{sec:llm}. All scans are parser-backed and reproducible.
No raw source code, file paths, or snippets are transmitted in research mode.

\textbf{Study~1 --- Governance System Audit.} A full deterministic scan of a 1,120-file
production Constitutional AI governance system developed with substantial AI coding assistance.
Governance-critical systems represent a high-stakes context for the failure mode \aira{}
targets. The developer was actively aware of and working to counter AI-induced failure
patterns throughout development.\footnote{The governance system is not named here to avoid
conflating the paper's findings with evaluation of that specific implementation. The pattern
is structural; the specific system is an instance of it.}

\textbf{Study~2 --- Corpus Pilot.} A rebuilt balanced comparison of 300 agent-attributed
files versus 300 matched human-control files, equally distributed across Python, JavaScript,
and TypeScript (100 files per language per arm). Agent-attributed files were sourced from
the AIDev dataset \citep{li2025aidev} (agentic PRs with confirmed \texttt{agent\_label}).
Human-control files were sourced from pre-2022 repositories with no AI tool references.
All scans were deterministic and parser-backed. The agent-attributed arm spans 126
repositories; the human-control arm spans 21 repositories.

\subsection{Study 1: Governance System Audit}
\label{sec:study1}

The scan produced 3,297 findings across 1,120 files: 1,283 \HIGH{}, 1,443 \MED{}, and
571 \LOW{}. Of 15 auditable checks, 13 returned \texttt{FAIL} and none returned
\texttt{PASS}. C07 and C12 returned \texttt{UNKNOWN} as expected.

The dominant signals were C03 (Broad Exception Suppression: 967 findings, 373 \HIGH{},
594 \MED{}) and C13 (Confidence Misrepresentation: 540 findings, all \MED{}). The
highest-density subsystem was the inference engine (408 files, 1,808 findings, 583 \HIGH{}),
followed by the test surface (544 files, 1,015 findings, 547 \HIGH{}). C14 (Test Coverage
Asymmetry) contributed 440 findings on the test surface alone --- the test suite itself
exhibits the happy-path bias the hypothesis predicts. This finding is directly comparable
to empirical results on test inclusion in agentic PRs, which show substantial variation in
test adoption rates across AI coding agents \citep{haque2026tests}.

A structurally notable property is \emph{severity clustering}: 11 of 13 automated checks
produce findings of exactly one severity level. C01, C02, C05, C09, C10, C11, and C15 are
exclusively \HIGH{}. C06, C08, and C13 are exclusively \MED{}. C04 is exclusively \LOW{}.
Only C03 and C14 span multiple levels. This reflects that different failure modes have
consistent consequence profiles --- checks that directly conceal failure produce only
\HIGH{} findings; checks that create ambiguity without direct concealment produce only \MED{}.

Four subsystems produced zero findings: the validation script, database layer, evidence
module, and uncertainty module. The framework discriminates rather than flagging
indiscriminately.

The density of findings in a codebase built by a developer actively working to counter
these patterns is a central result of Study~1. The failure distribution appears structural
rather than merely attentional.

\subsection{Study 2: Rebuilt Balanced Corpus Pilot}
\label{sec:study2}

\begin{table}[htbp]
\centering
\caption{Corpus Pilot --- Arm-Level Overview (300 files per arm)}
\label{tab:corpus-arm}
\small
\begin{tabular}{@{}llrrrr@{}}
\toprule
\textbf{Arm} & \textbf{Source} & \textbf{Files} & \textbf{HIGH} & \textbf{MED} &
\textbf{HIGH/file} \\
\midrule
A --- AI-attributed & AIDev (agentic PRs) & 300 & 80 & 148 & 0.267 \\
B --- Human baseline & Pre-2022 repos & 300 & 61 & 195 & 0.203 \\
\bottomrule
\end{tabular}
\end{table}

Table~\ref{tab:corpus-arm} summarizes the file-weighted comparison. The agent-attributed
arm produced 80 \HIGH{}, 148 \MED{}, and 24 \LOW{} findings, while the human-control arm
produced 61 \HIGH{}, 195 \MED{}, and 23 \LOW{}. This corresponds to 0.267 high-severity
findings per file versus 0.203, or a 1.32$\times$ excess in AI-attributed code.

\begin{table}[htbp]
\centering
\caption{High-Severity Findings per File by Language}
\label{tab:corpus-lang}
\footnotesize
\begin{tabularx}{\textwidth}{@{}>{\raggedright\arraybackslash}p{2.5cm}rrrr>{\raggedright\arraybackslash}X@{}}
\toprule
\textbf{Language} & \textbf{AI n} & \textbf{AI HIGH/f} &
\textbf{Human n} & \textbf{Human HIGH/f} & \textbf{Direction} \\
\midrule
JavaScript & 100 & 0.54 & 100 & 0.17 & AI $>$ Human (3.18$\times$) \\
Python & 100 & 0.11 & 100 & 0.12 & Parity \\
TypeScript (raw) & 100 & 0.15 & 100 & 0.32 & Human $>$ AI* \\
TypeScript (excl.\ outlier) & 100 & 0.15 & 99 & 0.10 & AI $>$ Human (1.49$\times$)* \\
\bottomrule
\end{tabularx}
\vspace{0.25em}
\parbox{\textwidth}{\footnotesize * See Section~\ref{sec:typescript-note} for repo-balance analysis.}
\end{table}

The separation was not uniform across languages (Table~\ref{tab:corpus-lang}). JavaScript
showed the clearest difference: 0.54 \HIGH{} per file in the agent-attributed arm versus
0.17 in the human-control arm. Python was approximately at parity (0.11 versus 0.12).

\subsubsection{TypeScript Repo-Balance Analysis}
\label{sec:typescript-note}

TypeScript initially appeared to cut against the hypothesis (0.15 versus 0.32 \HIGH{}/file).
However, follow-up analysis showed this counter-signal is not robust. The human-control
TypeScript arm spanned only 6 repositories, and a single file from \texttt{onlook-dev/onlook}
contributed 22 of the 32 human-control TypeScript \HIGH{} findings. Excluding that outlier
reduces the human-control TypeScript \HIGH{} rate from 0.32 to 0.101 per file.

A bootstrap sensitivity analysis (5,000 draws of 6 agent repositories from the 55 available)
found $P = 0.51$ for the agent mean equaling or exceeding the outlier-excluded human mean.
The apparent TypeScript reversal is not stable under repo-balanced analysis.

\begin{table}[htbp]
\centering
\caption{Check-Level Failure Counts by Arm (300 files per arm)}
\label{tab:corpus-checks}
\small
\begin{tabular}{@{}lrrl@{}}
\toprule
\textbf{Check} & \textbf{Arm A (AI)} & \textbf{Arm B (Human)} & \textbf{Direction} \\
\midrule
exception\_handling (C03) & 61 & 43 & AI $>$ Human \\
confidence\_representation (C13) & 11 & 7 & AI $>$ Human \\
fallback\_control (C04) & 14 & 12 & AI $>$ Human \\
\midrule
background\_tasks (C08) & 15 & 28 & Human $>$ AI \\
environment\_safety (C09) & 6 & 8 & Human $>$ AI \\
return\_contracts (C06) & 7 & 9 & Human $>$ AI \\
\bottomrule
\end{tabular}
\end{table}

At the check level (Table~\ref{tab:corpus-checks}), the strongest recurring support for the
hypothesis concentrated in exception-handling-related patterns. The agent-attributed arm
exceeded the human-control arm on \checkid{exception\_handling} (C03: 61 versus 43),
\checkid{confidence\_representation} (C13: 11 versus 7), and \checkid{fallback\_control}
(C04: 14 versus 12). By contrast, the human-control arm exceeded the agent-attributed arm
on \checkid{background\_tasks} (28 versus 15), \checkid{environment\_safety} (8 versus 6),
and \checkid{return\_contracts} (9 versus 7).

\subsection{LLM vs.\ Static Comparison}
\label{sec:llm}

To examine whether the failure-suppression pattern extends to the evaluation layer, we ran
a parallel hybrid scan using minimax-m2:cloud via Ollama as a secondary evaluator across the
Study~1 governance system surfaces. C07 and C12 were forcibly kept \texttt{UNKNOWN} in all
LLM output.

\begin{table}[htbp]
\centering
\caption{Static vs.\ LLM Finding Counts Across Governance System Surfaces}
\label{tab:llm}
\small
\begin{tabular}{@{}lrrr@{}}
\toprule
\textbf{Surface} & \textbf{Files} & \textbf{Static Findings} & \textbf{LLM Findings} \\
\midrule
Full repo (truncated) & 1,124 & 3,297 & 0 \\
engine/runtime & 28 & 222 & 5 \\
api & 58 & 250 & 6 \\
critic\_infrastructure.py & 1 & 34 & 6 \\
deliberation.py & 1 & 26 & 0 \\
startup.py & 1 & 29 & 0 \\
\bottomrule
\end{tabular}
\end{table}

\begin{table}[htbp]
\centering
\caption{Check-Level Suppression: Static FAIL vs.\ LLM PASS (Single Files, No Truncation)}
\label{tab:suppression}
\small
\begin{tabular}{@{}lrrrr@{}}
\toprule
\textbf{File} & \textbf{Static FAILs} & \textbf{LLM Suppressed} &
\textbf{Both FAIL} & \textbf{Suppression Rate} \\
\midrule
critic\_infrastructure.py & 7 & 2 & 5 & 29\% \\
deliberation.py & 7 & 7 & 0 & 100\% \\
startup.py & 8 & 8 & 0 & 100\% \\
\bottomrule
\end{tabular}
\end{table}

The single-file comparisons (Tables~\ref{tab:llm} and~\ref{tab:suppression}) are the most
controlled: no truncation, identical files, identical checks. For \texttt{deliberation.py},
the static scanner returned \texttt{FAIL} on 7 checks; the LLM returned \texttt{PASS} on
all 7 --- 100\% suppression. For \texttt{startup.py}, static \texttt{FAIL} on 8 checks;
LLM \texttt{PASS} on all 8 --- again 100\% suppression.

The suppression is not random across check types. The LLM consistently passed on C02, C03,
C04, C05, and C13 --- precisely the checks most directly related to failure concealment.
This result is not primarily a critique of the specific model used. It demonstrates that
the suppression tendency described in the Reward-Shaped Failure Hypothesis appears not only
in code generation but in AI-assisted code \emph{evaluation}. This motivates \aira{}'s
deterministic-first design: the audit instrument must not be susceptible to the failure mode
it is designed to detect.

\subsection{Discussion}

Study~1 shows that the predicted failure class is observable and measurable in a
production governance-critical codebase. The severity clustering result --- 11 of 13
automated checks producing single-severity findings --- suggests stable consequence profiles
across failure modes.

Study~2 sharpens rather than simplifies the hypothesis. The rebuilt balanced pilot supports
a real but conditional signal: AI-attributed code produces modestly more high-severity
findings overall (1.32$\times$), with the strongest separation in JavaScript (3.18$\times$)
and the clearest check-level support in exception-handling-related patterns. Python is near
parity. The apparent TypeScript reversal is substantially explained by control-side hotspot
concentration and does not survive repo-balanced sensitivity analysis.

The data therefore does not support the strongest version of the Reward-Shaped Failure
Hypothesis --- that AI-attributed code is uniformly hotter across all languages and surfaces.
It does support a more precise version: optimization pressure toward fail-soft behavior
produces the most consistent signal in JavaScript and exception-handling patterns. This
conditional account is scientifically stronger than a uniform claim, because it generates
testable predictions about where future studies should find the signal and where they
should not.

% ─────────────────────────────────────────────────────────────────────────────
\section{Output Specification}
\label{sec:output}
% ─────────────────────────────────────────────────────────────────────────────

All \aira{} audits produce output conforming to the following YAML structure (version~1.2):

\begin{lstlisting}[language=,caption={AIRA v1.2 Output Schema}]
ai_failure_audit:
  audit_version: "1.2"
  success_integrity:        PASS | FAIL | UNKNOWN
  audit_integrity:          PASS | FAIL | UNKNOWN
  exception_handling:       PASS | FAIL | UNKNOWN
  fallback_control:         PASS | FAIL | UNKNOWN
  bypass_controls:          PASS | FAIL | UNKNOWN
  return_contracts:         PASS | FAIL | UNKNOWN
  logic_consistency:        PASS | FAIL | UNKNOWN
  background_tasks:         PASS | FAIL | UNKNOWN
  environment_safety:       PASS | FAIL | UNKNOWN
  startup_integrity:        PASS | FAIL | UNKNOWN
  determinism:              PASS | FAIL | UNKNOWN
  lineage:                  PASS | FAIL | UNKNOWN
  confidence_opacity:       PASS | FAIL | UNKNOWN
  test_coverage_symmetry:   PASS | FAIL | UNKNOWN
  idempotency_safety:       PASS | FAIL | UNKNOWN
  findings:
    - issue: "<description>"
      file: "<path>"
      line: <number>
      severity: HIGH | MEDIUM | LOW
      check: "<check_id>"
\end{lstlisting}

\noindent\textbf{Data Availability.} The \aira{} scanner is released as open source at
\url{https://aira.bageltech.net} and \url{https://github.com/BDB-Labs/aira-scanner}.
Aggregate-only research submissions are supported via the CLI. Raw source code and file
paths are never transmitted in research mode. Benchmark datasets and additional comparative
scans are welcomed from the community.

% ─────────────────────────────────────────────────────────────────────────────
\section{Limitations}
\label{sec:limitations}
% ─────────────────────────────────────────────────────────────────────────────

\aira{} should be understood as a research scanner and evolving inspection framework.
The following limitations apply:

\begin{itemize}
  \item Cross-file and repo-level semantic reasoning remains weaker than single-file
        structural detection. Findings are primarily per-file and do not capture emergent
        failure patterns across module boundaries.

  \item The framework cannot determine that AI wrote the flagged code. The hypothesis
        predicts a higher rate of these patterns in AI-assisted codebases; \aira{} measures
        the patterns, not the authorship.

  \item A \texttt{PASS} result does not mean a system is safe. Absence of findings means
        the targeted patterns were not detected by the current rules --- not that the system
        exhibits full failure truthfulness.

  \item LLM-assisted scans at repo scale can become optimistic or lossy when prompts are
        truncated. The deterministic engine is the appropriate baseline for research use.

  \item Severity ratings are heuristic. Not every flagged pattern is a defect in context.
        Human review of findings is always required before remediation.

  \item Study~2 reports a rebuilt balanced pilot of 600 files. The pilot is file-balanced
        but not repo-balanced: the agent-attributed arm spans 126 repositories versus 21
        for the human-control arm. The TypeScript comparison is particularly sensitive to
        control-side hotspot concentration. Domain and system-type stratification are planned
        for the next corpus expansion.

  \item C07 and C12 remain human-review-only by design. Automated detection of parallel
        logic drift and source-to-output lineage failure requires repo-scale semantic
        comparison not yet implemented.
\end{itemize}

% ─────────────────────────────────────────────────────────────────────────────
\section{Conclusion and Future Work}
\label{sec:conclusion}
% ─────────────────────────────────────────────────────────────────────────────

This paper has introduced the Reward-Shaped Failure Hypothesis --- the proposal that AI coding
systems may systematically produce fail-soft code as a structural consequence of their
optimization environment --- and \aira{}, a 15-check inspection framework designed to detect
the specific failure class this hypothesis predicts.

The central reframing is this: traditional software validation asks whether systems work.
\aira{} asks whether systems tell the truth when they fail. Failure truthfulness is a
distinct and auditable system property. The 15 checks are organized around its measurement.

Study~1 results --- 3,297 findings across 1,120 files, 13 of 15 checks failing, and an LLM
evaluator under-reporting at a 44:1 ratio --- show that the predicted failure class is
observable, severe, and extends to the evaluation layer. Study~2, a rebuilt balanced 600-file
corpus pilot, provides early evidence beyond a single codebase: AI-attributed code shows a
1.32$\times$ overall excess of high-severity findings over matched human controls, with the
strongest signal in JavaScript (3.18$\times$) and exception-handling patterns. The hypothesis
is therefore supported as a real but conditional pattern --- not a uniform field-wide law,
but a measurable tendency whose strength varies by language and code surface.

The three most important next steps are: (1) corpus expansion with repo-balanced TypeScript
controls and domain/system-type stratification; (2) a controlled base-vs-RLHF model
comparison on structured generation tasks, to move the hypothesis from observational evidence
to direct experimental proof; and (3) expansion of the deterministic rule engine to support
cross-file and repo-level reasoning.

\aira{} is available as an open-source tool at \url{https://aira.bageltech.net} and
\url{https://github.com/BDB-Labs/aira-scanner}. The research collection pipeline supports
aggregate-only community submissions. Contributions are welcomed, especially comparative
scan datasets, rule extensions, and calibration studies.

% ─────────────────────────────────────────────────────────────────────────────
\section*{Acknowledgments}
% ─────────────────────────────────────────────────────────────────────────────

\aira{} was developed through direct observation during construction of Constitutional AI
governance frameworks with heavy AI coding assistance. The author thanks the broader AI
safety and software reliability research communities whose parallel work motivates the bridge
this paper attempts to build.

% ─────────────────────────────────────────────────────────────────────────────
\bibliographystyle{plainnat}
\begin{thebibliography}{9}

\bibitem[Beyer et~al.(2016)]{beyer2016sre}
B.~Beyer, C.~Jones, J.~Petoff, and N.~Murphy.
\newblock \emph{Site Reliability Engineering: How Google Runs Production Systems}.
\newblock O'Reilly Media, 2016.

\bibitem[Casper et~al.(2023)]{casper2023open}
S.~Casper, X.~Davies, C.~Shi, T.~Gilbert, J.~Scheurer, J.~Rando, R.~Freedman,
T.~Korbak, D.~Lindner, P.~Freire, et~al.
\newblock Open problems and fundamental limitations of reinforcement learning from
human feedback.
\newblock \emph{arXiv preprint arXiv:2307.15217}, 2023.

\bibitem[Gao et~al.(2025)]{gao2025survey}
R.~Gao et~al.
\newblock A survey of bugs in AI-generated code.
\newblock \emph{arXiv preprint arXiv:2512.05239}, 2025.

\bibitem[Haque et~al.(2026)]{haque2026tests}
S.~Haque, D.~Str{\"u}ber, and N.~Tsantalis.
\newblock Do autonomous agents contribute test code? {A} study of tests in
agentic pull requests.
\newblock \emph{arXiv preprint arXiv:2601.03556}, 2026.

\bibitem[Li et~al.(2025)]{li2025aidev}
H.~Li, H.~Zhang, et~al.
\newblock {AIDev}: The rise of {AI} teammates in software engineering~3.0.
\newblock Dataset available at \url{https://huggingface.co/datasets/hao-li/AIDev}, 2025.

\bibitem[MSR(2026)]{msr2026testfailures}
MSR 2026 Mining Challenge.
\newblock An empirical analysis of test failures in AI-generated pull requests.
\newblock In \emph{Proc.\ 23rd International Conference on Mining Software Repositories
(MSR)}, Rio de Janeiro, Brazil, April 2026.

\bibitem[Ogenrwot and Businge(2026)]{ogenrwot2026how}
E.~Ogenrwot and J.~Businge.
\newblock How {AI} coding agents modify code: {A} large-scale study of {GitHub}
pull requests.
\newblock \emph{arXiv preprint arXiv:2601.17581}, 2026.

\bibitem[Yu et~al.(2024)]{yu2024codereval}
H.~Yu, W.~Shen, K.~Ran, J.~Liu, Q.~Wang, and Y.~Jiang.
\newblock {CoderEval}: A benchmark of pragmatic code generation with generative
pre-trained models.
\newblock In \emph{Proc.\ IEEE/ACM ICSE}, 2024.

\end{thebibliography}

\end{document}
